\section{Local conventions and parameter theorems}\label{sec:parameters}

We record the precise parameter theorems used below and explain how their
semisimple Weil parameters relate to monodromy. The distinction is essential:
the global theorem and the tempered-completion theorem will produce two
potentially different monodromy operators over the same Weil parameter.

\subsection{Groups, induction, and Whittaker characters}

Let $k$ be a nonarchimedean local field of positive characteristic, with
residue field of cardinality $Q$. Fix a split adjoint absolutely simple
group $G$ over $k$, a split maximal torus $T$, and a Borel subgroup $B=TU$.
In the local arguments, $H$ denotes either $G$ or a standard Levi subgroup
of $G$; a standard parabolic of $H$ is written $P=MR$. Dual groups have
compatible positive systems, so positive coroots of $H$ are positive roots
of $\widehat H$. Put $\widehat{\mathfrak n}=\Lie(\widehat R)$.
All parabolic induction is normalized and is denoted $I_P^H$.

For a split group, a Whittaker character is a smooth unitary character of
its maximal unipotent subgroup that is nontrivial on every simple root
group. A smooth representation is \emph{generic} if it admits a nonzero
equivariant functional for such a character. We use uniqueness of Whittaker
functionals for irreducible generic representations, heredity under
normalized induction, and exactness of twisted unipotent coinvariants.
These statements, and the local-coefficient constructions based on them,
are valid over local function fields; see \cite[Sections~1--2]{LomeliLS}
and \cite[Section~2]{DCHL}. Exactness is also
\cite[Proposition~1.9(a)]{BernsteinZelevinsky}. For the foundational uniqueness and heredity
results, see also \cite[Th\'eor\`emes~2 and~4]{Rodier}.

The following elementary observation ensures that Whittaker data can be
chosen compatibly throughout the argument, even in small characteristic.

\begin{lemma}\label{lem:whittaker}
Every standard Levi subgroup $H$ of $G$ has split connected center.
Every smooth character of its maximal unipotent subgroup is trivial on
the nonsimple positive root groups. Its nondegenerate unitary characters
form a single $T(k)$-orbit.
\end{lemma}

\begin{proof}
Because $G$ is adjoint, its simple roots are a $\mathbb Z$-basis of
$X^*(T)$. The center of $H$ is the simultaneous kernel of a subset of
these coordinate characters, hence is a split torus. The same basis
allows arbitrary independent rescaling of the simple-root coordinates
belonging to $H$. Once triviality on nonsimple root groups is known,
additive self-duality of $k$ proves the orbit assertion.

To prove that triviality, a nonsimple positive root may be written as a
sum of two positive roots. The full root subsystem in their span reduces
the question to $A_2$, $B_2$, or $G_2$, with its induced positive system.
In $A_2$ the usual commutator is a parametrization of the nonsimple root
group. For $B_2$ or $G_2$, write $a,b$ for the short and long simple roots.
In $G_2$ first kill the highest root group $U_{3a+2b}$ using
$[U_b,U_{3a+b}]$, whose coefficient is $\pm1$. Modulo that group, the
commutator $[x_a(u),x_b(t)]$ has factors
\[
 x_{ja+b}(\epsilon_j u^j t),\qquad \epsilon_j\in\{1,-1\},
\]
with $1\leq j\leq2$ in $B_2$ and $1\leq j\leq3$ in $G_2$.
Fix a nontrivial additive character $\psi$ of $k$, and write the
restriction of the given character to $U_{ja+b}$ as
$x_{ja+b}(t)\mapsto\psi(c_jt)$. It kills commutators, so
\[
 \psi\left(t\sum_j\epsilon_jc_ju^j\right)=1
 \quad\text{for every }u,t\in k.
\]
Varying $t$ gives $\sum_j\epsilon_jc_ju^j=0$ for every $u$.
Since $k$ is infinite, all $c_j$ vanish. This polynomial argument uses no
division by $2$ or $3$ and remains valid in those characteristics.
\end{proof}

Write $\mathfrak a_{M,\C}^*=X^*(M)\otimes_{\mathbb Z}\C$.
For $z\in\mathfrak a_{M,\C}^*$, the notation $\sigma_z$ denotes the
unramified twist in which $s\chi$ acts as $|\chi|^s$.
Via duality, $X^*(M)$ is the cocharacter lattice of
$Z(\widehat M)^\circ$; a central weight $b$ on a representation of
$\widehat M$ therefore pairs with $z$, giving $b(z)$.
In particular, the weights on $\widehat{\mathfrak n}$ satisfy
$b(\nu)>0$ when $\nu$ is in the open positive chamber for $P$.

\subsection{Weil--Deligne parameters and local factors}

Fix $\ell\ne\operatorname{char}(k)$ and an isomorphism
$\iota:\overline{\Q}_\ell\simeq\C$. All absolute values of coefficients
below are taken after $\iota$. Choose the square roots in normalized
induction and in the parameter theorems to correspond to positive real
square roots. We use geometric Frobenius $\Fr$, so $|\Fr|=Q^{-1}$.
Local class field theory sends a uniformizer to geometric Frobenius,
and Satake parameters are normalized accordingly. Thus a twist by
$|\chi|^s$ on representations gives the same $|\cdot|^s$-twist on
parameters.

A Frobenius-semisimple Weil--Deligne parameter for $H$ is a pair $(r,N)$,
where $r:W_k\longrightarrow\widehat H(\C)$ has finite inertia image and
semisimple Frobenius, and $N\in\Lie(\widehat H)$ is nilpotent with
\begin{equation}\label{eq:wd-relation}
 \Ad(r(w))N=|w|N.
\end{equation}
For an algebraic representation of $\widehat H$ on $V$, use the same
symbols for the induced operators and set
\begin{equation}\label{eq:local-factor}
 L(s,V)=\det\left(1-Q^{-s}r(\Fr)\mid(\ker N)^{r(I_k)}\right)^{-1}.
\end{equation}
This is a reciprocal polynomial in $Q^{-s}$, with constant term $1$;
in particular, it has no zeros. We often write $L(s,a\circ(r,N))$ when
the algebraic representation $a$ needs to be specified.

A tempered $L$-parameter is a homomorphism
$\phi:W_k\times\SL_2(\C)\longrightarrow\widehat H(\C)$, algebraic on
$\SL_2$, whose Weil image is bounded. Its associated pair is
\begin{equation}\label{eq:completion}
 r(w)=\phi\left(w,
       \begin{pmatrix}|w|^{1/2}&0\\0&|w|^{-1/2}\end{pmatrix}\right),
 \qquad
 N=d\phi\begin{pmatrix}0&1\\0&0\end{pmatrix}.
\end{equation}
Thus the Weil part $r$ of a tempered completion need not itself be
bounded. This familiar distinction is the reason to retain monodromy.

We write $f\doteq g$ for equality up to a nonzero constant times a
monomial in the exponential unramified-twist coordinates. Finite covers
of twist tori are allowed. In additive coordinates, the omitted factor
is a constant times the exponential of a linear form, hence is
holomorphic and nowhere zero. In one variable it has the form $cQ^{-as}$.

\subsection{The parameter inputs}

We use the following forms of the global and local parameter theorems.
They specify the precise information carried from representations to
parameters; no compatibility of monodromy is included.

\begin{theorem}[Lafforgue; Genestier--Lafforgue]\label{thm:parameters}
For the split groups considered here the following hold.
\begin{enumerate}[label=\textup{(\roman*)}]
\item An irreducible smooth local representation $\sigma$ has a
semisimple Weil parameter $\rho_\sigma$, with finite inertia image.
The parametrization is compatible with normalized parabolic induction,
tori and local class field theory, group isomorphisms, central characters,
and twists through algebraic characters of the group.
\item A cuspidal automorphic representation $\Pi$ over a function field,
with finite-order central character, occurs in an excursion summand
with a semisimple global parameter $\Sigma_\Pi$. This parameter is
defined over a finite extension of $\Q_\ell$ and matches normalized
Satake parameters away from a finite set.
\item At every place $v$, the semisimplification of
$\Sigma_\Pi|_{W_{F_v}}$ is conjugate to $\rho_{\Pi_v}$.
\end{enumerate}
\end{theorem}

The global statement is \cite[Th\'eor\`eme~0.1]{VLafforgue}; the local
statements, including representations without an integrality condition,
are \cite[Th\'eor\`eme~0.1 and Remarque~0.2]{GenestierLafforgue}.
For compatibility with central characters, apply the functoriality for
homomorphisms with normal image to the inclusion of the central split
torus.

Here is how to interpret the coefficient and occurrence assertions for
complex representations. A finite-order central character allows a
cocompact central lattice in its kernel. Fixed-level cusp spaces modulo
that lattice are finite-dimensional and commute with coefficient
extension. Project an irreducible Hecke module of level invariants to
any excursion summand where its projection is nonzero; that projection
is injective, so it supplies the required occurrence. Uniqueness of the
excursion parameter of $\Pi$ is unnecessary. For the adjoint group no
central lattice is needed.

Locally, after transport by $\iota$, irreducibles may be realized over a
finite extension of $\Q_\ell$. Indeed a fixed-compact-open Hecke algebra
over $\Q_\ell$ is of finite type
\cite[Theorem~1.1]{DHKMFiniteness2024}, so its finite-dimensional simple
module over $\overline{\Q}_\ell$ descends to
such an extension. The corresponding irreducible representation also
descends: form induction from this Hecke module and quotient by the
largest subrepresentation with zero invariants under that compact
open subgroup. The latter subrepresentation is characterized by
vanishing of every averaged translate, a condition compatible with
scalar extension. This permits use of the finite-extension formulations
of the parameter theorems.

\begin{theorem}[Gan--Harris--Sawin, with Beuzart-Plessis]
\label{thm:tempered-completion}
If $\sigma$ is tempered on a standard Levi subgroup $M$ of $G$, then
$\rho_\sigma$ has a tempered completion $(\rho_\sigma,N_M)$ as in
\eqref{eq:completion}. This holds in every positive characteristic.
\end{theorem}

\begin{proof}[Explanation of the cited form]
The published \cite[Theorem~1.2 and Corollary~1.3]{GHS} give an
\emph{essentially} tempered completion for arbitrary connected reductive
groups over local function fields. Thus its Weil image is bounded modulo
the center. Since $\sigma$ is tempered, its central character is unitary.
By Theorem~\ref{thm:parameters}, the dual morphism
$\widehat M\to\widehat{Z(M)}$ has bounded image on the parameter.
Root data show that its characters span the rational character space
of the abelian quotient of $\widehat M$. Moreover, the $\SL_2$ factor
dies in that quotient. Hence the Weil image of the completion is
bounded in the abelian quotient as well. The map from $\widehat M$ to
the product of its adjoint and abelian quotients has finite kernel;
boundedness in these two quotients therefore gives boundedness in
$\widehat M$ itself.
\end{proof}

\subsection{The Weil part of a global restriction}

The local monodromy theorem turns a restriction of $\Sigma_\Pi$ into a
Weil--Deligne pair. The next observation verifies that its
Frobenius-semisimple Weil part is exactly the parameter in
Theorem~\ref{thm:parameters}(iii), also as a dual-group-valued parameter.

\begin{lemma}\label{lem:wd-semisimplification}
Let $\Sigma$ be a continuous $\ell$-adic representation of a local
decomposition group into $\widehat H$, defined over a finite extension
of $\Q_\ell$. If $(r,N)$ is its Frobenius-semisimplified Weil--Deligne
pair, then $r$ is conjugate to the semisimplification of $\Sigma|_{W_k}$.
\end{lemma}

\begin{proof}
Before Frobenius semisimplification, write the pair as $(r_0,N)$, with
$r_0(I_k)$ finite. The construction takes place in $\widehat H$:
in a faithful linear representation the unipotent logarithm lies in
its Lie algebra, and $\Sigma$ on the Weil group is obtained from $r_0$
by multiplying by the appropriate exponentials of $N$.

Write $r_0(\Fr)=hu$ for its commuting semisimple and unipotent parts.
Some positive power of $r_0(\Fr)$ centralizes the finite inertia image.
Jordan decomposition in that centralizer shows that a power of $u$,
and hence $u$, centralizes it. The relation
$\Ad(hu)N=Q^{-1}N$ gives $\Ad(h)N=Q^{-1}N$ and $\Ad(u)N=N$.
Replacing $hu$ by $h$ thus defines a Weil parameter $r$ with the same
finite inertia. The identity component of the Zariski closure of $r(W_k)$
is toral, so $r$ is semisimple.

Let $S$ be the identity component of the Zariski closure of the cyclic
group generated by $h$. It is a torus centralizing $r_0(I_k)$ and $u$.
If $N\ne0$, its weight on the line $\C N$ is nontrivial, because
$Q^{-1}$ is not a root of unity. A cocharacter of $S$ can therefore
contract $N$ to zero while fixing $r_0$. This realizes removal of the
monodromy exponentials as a conjugation limit in $\widehat H$.
Next, the centralizer of $r$ has reductive identity component and contains
the unipotent element $u$. A cocharacter in that centralizer contracts
$u$ to the identity. These two limits preserve semisimplification and
leave precisely $r$. Finally, density of the Weil group in the
decomposition group gives the same Zariski closure for both images.
\end{proof}

Thus a global restriction and a tempered local completion can have the
same $r$ while retaining distinct operators $N$. The next section gives
two ways to work with this distinction: an $L$-factor ratio that forgets
$N$, and a regularity condition that determines $N$ up to conjugacy.
